% !TeX program = xelatex
\documentclass[preprint,review,12pt,a4paper]{elsarticle}

% ── 中文支持（XeLaTeX）─ scheme=plain 避免与 elsarticle frontmatter 冲突 ──
\usepackage[scheme=plain]{ctex}
\setmainfont{Times New Roman}
\setCJKmainfont{SimSun}[BoldFont=SimHei,ItalicFont=KaiTi]

% ── 页面 ── (preprint 模式下 elsarticle 不加载 geometry，需手动设)
\usepackage[top=2.5cm,bottom=2.5cm,left=2.5cm,right=2.5cm]{geometry}

% ── 图表 ──
\usepackage{graphicx}
\usepackage{caption}
\usepackage{subcaption}
\usepackage{float}
\usepackage{booktabs}
\usepackage{multirow}
\graphicspath{{figures/}}

% ── 表格 ──
\usepackage{array}
\usepackage{longtable}
\usepackage{makecell}

% ── 数学 ──
\usepackage{amsmath}
\usepackage{amssymb}

% ── 超链接 ──
\usepackage[unicode=true]{hyperref}
\hypersetup{colorlinks=true,citecolor=blue,linkcolor=black,urlcolor=blue}

% ── 汉化标签 ──
\abstracttitle{摘要}
\keywordtitle{关键词}
\renewcommand{\appendixname}{附录}

% ── 图注字号 ── (覆盖 elsarticle 的 \footnotesize，中文太小)
\captionsetup{font=small}

\begin{document}

\begin{frontmatter}

\title{基于多路段自然驾驶轨迹的变道风险跨路段泛化预测与多模型SHAP可解释性分析}

\author{樊琼}
\author{沈照庆\corref{cor1}}
\cortext[cor1]{沈照庆（1981—），男，博士，长安大学公路学院，副教授，主要研究方向为道路交通工程。}
\affiliation{organization={长安大学公路学院}, city={西安}, country={中国}}

\begin{abstract}
针对变道风险预测模型在路段间可迁移性不明确这一关键问题，本研究基于 CQSkyEyeX 自然驾驶轨迹数据集，系统评估了四种机器学习模型在 5 个高速公路场景（路基段、桥梁段、分合流区）中的跨路段泛化性能。首先，通过整合 mTTC、THW、PET 和 ETTC 等多维度替代安全指标，构建了速度自适应的加权评分体系，实现变道风险的三级标注（高风险 378 辆、中风险 561 辆、低风险 635 辆，共计 1,574 个变道样本）。随后，对 XGBoost、Random Forest、MLP 和 LSTM 四种模型分别进行留一路段交叉验证与随机 80/20 划分对比评估。实验结果表明：树集成模型在跨路段泛化中表现显著优于神经网络模型 —— XGBoost 取得最优跨路段性能（Accuracy=0.764$\pm$0.078, Weighted F1=0.764$\pm$0.079），领先 Random Forest 约 3.2 个百分点，而 MLP 和 LSTM 在面对复杂跨路段场景时泛化能力明显不足。路段几何形态对模型泛化能力具有显著影响：路基段间可迁移性最高（Fold 2 F1=0.875），分合流区最难以泛化（Fold 5 F1=0.630）。SHAP 分析一致确认车头时距最小值（Time\_Headway\_min）为跨路段最稳定的风险判据（XGBoost 高风险贡献值 +1.251），后方距离与跟车距离的重要性则呈现路段间差异。本研究系统地基于中国高速公路多场景数据，探讨了变道风险模型跨路段泛化性，为实际变道预警系统的模型选择与路段自适应部署提供了实证依据和工程参考。
\end{abstract}

\begin{keyword}
变道风险评估 \sep 机器学习 \sep 跨路段泛化 \sep SHAP 可解释性 \sep 替代安全指标 \sep 自然驾驶轨迹
\end{keyword}

\end{frontmatter}

% ── 1 引言 ──
\section{引言}

变道操作（Lane Change, LC）相关碰撞占总事故的 4\%$\sim$10\%\cite{C01,C02}，其风险评定需综合纵向碰撞与横向多车交互\cite{C03}。因历史碰撞数据稀疏，研究者引入 TTC\cite{C05}、PET\cite{C04}、SDI\cite{C06}等替代安全指标（SSMs）以量化潜在冲突。然而传统 SSMs 在变道中存在双重局限：多数指标仅适用跟驰场景，无法刻画横向交互时变性\cite{C03,C07}；单维度难以覆盖变道事件的多风险源特性\cite{C01}。为此研究者发展了多种综合评估框架——Park等\cite{C08}通过故障树分析提出变道风险指数（LCRI），Wu等\cite{C03}构建了时空风险估计（TSRE）模型，Xue等\cite{C02}引入安全裕度概念。近年来ML方法被引入变道风险预测领域，冯移冬等\cite{C09}利用RF-RFE特征选择实现汇入场景F1=84\%，温惠英等\cite{C10}使用XGBoost+SHAP揭示合流区致险因子；DS-LSTM\cite{C11}和CNN-BiLSTM-Attention\cite{C12}则展示了深度学习在时序建模中的潜力。SHAP\cite{C13}和TreeSHAP\cite{C14}的提出为模型可解释性提供了统一框架。

然而一个根本性问题尚未系统回答：\textbf{在同一场景训练的变道风险模型能否有效迁移至几何形态或限速不同的其他路段？}现有研究多局限于单一数据集\cite{C15}，少数跨数据集证据表明风险因子存在站点间异质性\cite{C16}，但系统性的多路段多模型对比评估仍然缺乏。该问题的答案直接决定了变道预警系统是否需要为每个路段重新训练专用模型，具有显著的工程与经济含义。

\subsection{相关研究简述}

替代安全指标（SSMs）自 Hayward\cite{C05}提出 TTC 以来，已发展出涵盖时间型、距离型和减速度型的完整体系\cite{C18}。TTC 在变道动态多车交互中频繁失效\cite{C01}，Minderhoud 与 Bovy\cite{C06}发展出 TET 和 TIT 等扩展指标。在综合评估方面，Park 等\cite{C08}构建了变道风险指数（LCRI），Wu 等\cite{C03}发展了 TSRE 模型，Xue 等\cite{C02}引入安全裕度概念。风险场理论方面，刘振宇与项巧军\cite{C01}构建 F-F Diagram，王涛等\cite{C07}实现了从确定性到概率性评估的范式转变，胡立伟等\cite{C29}进一步提出了考虑风险脉冲效应的统一场感知模型。余加勇等\cite{C26}基于故障树法融合时空指标构建了换道动态风险综合量化方法。

ML 方法在变道风险预测中展现出优于传统方法的性能。齐龙等\cite{C30}利用K-means++聚类与CNN-LSTM实现了多车道合流区换道识别与轨迹预测。邢璐等\cite{C31}考虑客货交互压迫度构建两阶段换道风险预测模型。冯移冬等\cite{C09}利用 RF-RFE 特征选择实现了汇入场景 F1=84\%；温惠英等\cite{C10}使用 XGBoost+SHAP 揭示合流区风险因子。树集成方法——XGBoost\cite{C19}和 RF\cite{C20}——因对表格化数据的高度适应性而被广泛采用。深度学习中，DS-LSTM\cite{C11}融合驾驶风格聚类，CNN-BiLSTM-Attention\cite{C12}实现多车型变道风险识别。郭孜政等\cite{C28}基于领域自适应方法实现了驾驶人分心行为检测。SHAP\cite{C13} 和 TreeSHAP\cite{C14} 为模型可解释性提供了理论框架，但现有研究多停留于全局特征重要性排序，未触及跨场景维度。

在跨路段泛化方向，Huang 等\cite{C22}发现模型在桥隧段训练后部署至合流段时可靠度显著下降；Wang 等\cite{C16}确认了风险因子的站点间异质性。夏萧菡等\cite{C32}利用异常驾驶行为数据实现了高速公路实时事故风险预测，魏伟等\cite{C34}构建了基于深度学习的交通冲突预测方法。Ahmed 等\cite{C23}的域自适应策略提升有限（3\%$\sim$5\%）。

现有研究存在三个局限：（1）多使用仿真场景，缺乏实际复杂道路覆盖；（2）缺少树模型与神经网络泛化行为的系统比较；（3）尚未结合 SHAP 从特征重要性动态变化角度解释泛化退化内因。针对上述不足，本文基于 CQSkyEyeX 数据集\cite{C17}中 5 种不同道路几何形态的实际场景，开展变道风险 ML 预测模型的跨路段泛化性系统研究。主要贡献如下：

\begin{enumerate}
  \item \textbf{多路段评估基准构建}：基于涵盖路基段、整体式桥梁段和分合流区三种道路类型的 1,574 个变道样本，建立了多指标加权安全评分 $\to$ 三级风险标签（高风险 378 辆、中风险 561 辆、低风险 635 辆）的完整标注管线，为跨路段泛化评估提供了统一的数据基础。
  \item \textbf{四模型跨路段泛化对比}：通过留一路段交叉验证（5-fold cross-location CV）与随机划分的对照设计，量化了 XGBoost、Random Forest、MLP 和 LSTM 四种代表性模型的跨路段性能退化程度，揭示了树集成模型与神经网络模型在泛化行为上的本质差异。
  \item \textbf{SHAP 驱动的风险因子普适性分析}：利用 XGBoost 和 Random Forest 的 SHAP 值多模型一致性分析，首次在跨路段维度上分离了变道风险的普适性特征（通用因子）与路段特异性特征（场景因子），为路段自适应模型设计提供了特征层面的依据。
\end{enumerate}

\section{变道风险预测模型构建}

\subsection{CQSkyEyeX 自然驾驶轨迹数据来源与预处理}

本研究采用 CQSkyEyeX 自然驾驶轨迹数据集\cite{C17}\footnote{数据集官网：http://www.cqskyeyex.com/CQSkyEyeX}。该数据集基于无人机航拍视频经计算机视觉方法提取车辆轨迹\cite{C24}，定位精度 RMSE $<0.10$ m。原始采样率 30 Hz，研究中统一降采样至 5 Hz。

本研究选取数据集中已被完整标注和公开的前 5 个场景，其道路几何形态涵盖三种典型高速公路段落类型：路基段、整体式桥梁段和分/合流区。表~\ref{tab:scenes} 汇总各场景的详细特征。

\begin{table}[H]
\centering
\caption{CQSkyEyeX 数据集 5 场景特征描述}
\label{tab:scenes}
\small
\begin{tabular}{lcccccc}
\toprule
场景 &  路段类型 & 单向车道数 & 限速 (km/h) & 拍摄长度 (m) & 变道样本数 \\
\midrule
Loc1 & 路基基本段 & \multirow{5}{*}{3} & \multirow{4}{*}{120} & \multirow{4}{*}{$\!\sim\!420$} & 151 \\
Loc2 & 路基基本段 &                    &                      &                               & 243 \\
Loc3 & 整体式桥梁段 &                  &                      &                               & 440 \\
Loc4 & 整体式桥梁段 &                  &                      &                               & 306 \\
\cmidrule(lr){4-5}
Loc5 & 分/合流区 &                    & 100                  & $\!\sim\!350$                 & 434 \\
\bottomrule
\end{tabular}
\end{table}

图~\ref{fig:velocity_distribution} 展示了 3 个代表性场景的速度分布密度。Loc1（路基段）呈明显双峰结构（货车低速峰与客车高速峰），Loc3（桥梁段）双峰弱化，Loc5（分合流区）趋于单峰集中。王博等\cite{C27}关于货车运行速度特性的研究也佐证了路段几何形态对速度分布的差异化影响。这种分布差异决定了 $K$ 值对幂指数 $B$ 的敏感程度：Loc5 中低速集中使 $K\equiv1$，而 Loc1 高速峰群体中 $K$ 可产生有效修正。

\begin{figure}[H]
\centering
\begin{subfigure}{0.32\textwidth}
\centering
\includegraphics[width=\textwidth]{velocity_following_distribution_location1.png}
\caption{Loc1（路基段）}
\end{subfigure}
\begin{subfigure}{0.32\textwidth}
\centering
\includegraphics[width=\textwidth]{velocity_following_distribution_location3.png}
\caption{Loc3（桥梁段）}
\end{subfigure}
\begin{subfigure}{0.32\textwidth}
\centering
\includegraphics[width=\textwidth]{velocity_following_distribution_location5.png}
\caption{Loc5（分合流区）}
\end{subfigure}
\caption{三个代表性场景跟驰车辆速度分布密度曲线。Loc1 呈现明显车型双峰（货车低速峰 vs 小客车高速峰）；Loc3 双峰弱化；Loc5 趋于单峰集中。}
\label{fig:velocity_distribution}
\end{figure}

原始数据集提供邻车 ID，各方向邻车的空间间距由坐标几何关系计算得到。变道样本以车道线跨越时刻 $t_0$ 为参考，截取 $t_{-100}$ 至 $t_{-25}$ 共 75 帧（3 s\cite{C02,C03}），经移动平滑去噪后获得 1,574 个有效样本。变道窗口内各风险指标具有时序非平稳性，支持采用均值/标准差/最值等聚合特征捕获极端状态信息而非单帧瞬时值的特征工程策略。

\subsection{基于熵权法的替代安全指标与危险评分体系构建}

为实现对变道轨迹的无监督风险标注，同时克服传统离散阈值打分法中类别边界僵硬、对指标小幅度变化不敏感的缺陷，本文采用基于指数衰减函数的连续风险评分方案。评分中的权重 $\{w_i\}$ 由熵权法（Entropy Weight Method, EWM）从数据中客观确定。EWM 是一种无需标签的客观赋权方法，其基本逻辑为：信息越分散（各车辆在该指标上的激活值差异越大）的指标越具判别力，权重越高。具体计算流程如下：

\textbf{Step 1 — 风险激活矩阵。} 对 $n$ 辆变道车辆和 $m=5$ 项 SSM 指标（mTTC, THW, PET, F\_ETTC, OL\_PET），逐车取全窗口最危险帧的指数衰减值构建 $\mathbf{X}\in[0,1]^{n\times m}$：
\begin{equation}
X_{ij} = \max_{t \in [1,\,75]}\; \exp\!\left(-\frac{v_j^{(i)}(t)}{k_j}\right)
\end{equation}
若第 $i$ 辆车不存在第 $j$ 个指标对应的邻车，则 $X_{ij}=0$。

\textbf{Step 2 — 比重矩阵。} 列归一化得 $\mathbf{P}$，其中 $P_{ij}=X_{ij}/\sum_{i=1}^n X_{ij}$：
\begin{equation}
P_{ij} = \frac{X_{ij}}{\sum_{i=1}^{n} X_{ij}}
\end{equation}

\textbf{Step 3 — 信息熵与差异系数。} 第 $j$ 个指标的信息熵 $E_j=-\frac{1}{\ln n}\sum_i P_{ij}\ln P_{ij}$，差异系数 $D_j=1-E_j$。$E_j$ 越大（信息越集中）则权重越低。

\textbf{Step 4 — 权重归一化。} $w_j=D_j/\sum_k D_k$，满足 $\sum w_j=1$。

EWM 权重仅依赖 5 项 SSM 指标的风险激活值矩阵 $\mathbf{X}$，完全不访问模型训练所用的 60 维样本特征和风险标签，保证了权重标定的统计独立性和可复现性。由上述流程得出的变道车辆权重为 $\{w_j\}^{5}_{j=1} = \{0.247, 0.083, 0.131, 0.356, 0.183\}$（mTTC/THW/PET/F\_ETTC\newline/OL\_PET），跟驰车辆权重为 $\{0.463, 0.114, 0.423\}$（mTTC/THW/B\_mTTC）。值得关注的是 F\_ETTC 获得了最高权重（0.356），显著高于其经验估计值（0.20），表明前向扩展碰撞时间在变道车辆间具有最大的判别力差异——即不同变道事件中前向碰撞风险的异质性远超其他四个指标；而 THW 仅获 0.083，说明车头时距的分布在车辆间高度集中（$E_j$ 接近 1），提供的额外判别信息有限。

\subsubsection{变道场景的车辆交互关系}

在变道执行过程中，自车（Subject Vehicle, SV）与周边四辆邻车构成典型的五车交互系统（图~\ref{fig:lc_interaction}）：当前车道正前方车辆（Preceding Vehicle, PV）、当前车道正后方车辆（Following Vehicle, FV）、目标车道正前方车辆（Lead Vehicle, LV）、目标车道正后方车辆（Lag Vehicle, LgV）。不同的 SSM 指标度量的正是 SV 与特定邻车之间的冲突紧迫度，明确这一映射关系是后续风险评分公式的物理前提。

\begin{figure}[H]
\centering
\includegraphics[width=0.6\textwidth]{fig_lc_interaction.png}
\caption{变道车辆（SV）与周边四辆邻车的交互关系。LV：目标车道前车（Lead Vehicle）；LgV：目标车道后车（Lag Vehicle）；PV：当前车道前车（Preceding Vehicle）；FV：当前车道后车（Following Vehicle）。箭头示意车辆行驶方向与横向偏移意图。}
\label{fig:lc_interaction}
\end{figure}

本节所定义的五个替代安全指标与上述车辆对的对应关系如下：mTTC 和 THW 度量 SV 与 PV 之间的纵向碰撞风险（当前车道追尾）；PET 度量 SV 与目标车道邻车（LV 或 LgV）先后通过同一空间点的时间逼近程度；F\_ETTC 将碰撞时间拓展至二维平面，度量 SV 与 LV 沿碰撞危险方向的综合逼近速率；OL\_PET 度量 SV 横向侵入目标车道时与 LgV 的侧向时间裕度。跟驰风险评分（式 8）中的 B\_mTTC 则度量 FV 与 SV 之间的后向追尾风险。

\subsubsection{替代安全指标的数学定义}

\textbf{(1) 修正碰撞时间（mTTC）。} 传统 TTC 假设运动状态恒定而忽略加速度，mTTC将加速度纳入碰撞时间估计。设两车相对位置矢量为 $\Delta\mathbf{s} = (\Delta x, \Delta y)$，相对速度矢量为 $\Delta\mathbf{v} = (\Delta v_x, \Delta v_y)$，相对加速度矢量为 $\Delta\mathbf{a} = (a_x, a_y)$，则：

\begin{equation}
\text{mTTC} = \frac{-\langle\Delta\mathbf{v}, \Delta\mathbf{s}\rangle + \sqrt{\langle\Delta\mathbf{v}, \Delta\mathbf{s}\rangle^2 + 2\langle\Delta\mathbf{a}, \Delta\mathbf{s}\rangle\|\Delta\mathbf{s}\|^2}}{2\langle\Delta\mathbf{a}, \Delta\mathbf{s}\rangle}
\end{equation}

其中 $\langle\cdot,\cdot\rangle$ 表示向量的内积。当后车速度不大于前车（分母为 0 或根号内为负）时，$\text{mTTC} \to \infty$，表示当前运动状态下不可能发生碰撞。

\textbf{(2) 车头时距（THW）。} 表征自车以当前速度到达前车当前位置所需的时间：

\begin{equation}
\text{THW} = \frac{d_{\!f}}{v_e}
\end{equation}

其中 $d_{\!f}$ 为自车前端至前车尾端的纵向间距（m），$v_e$ 为自车速度（m/s）。THW 直接反映了驾驶人可用的反应时间储备。

\textbf{(3) 侵入后时间（PET）。} 度量两车轨迹在空间中先后通过同一冲突点的时间差\cite{C04}：

\begin{equation}
\text{PET} = t_{2} - t_{1}
\end{equation}

其中 $t_1$ 为第一辆车离开冲突区域的时刻，$t_2$ 为第二辆车到达冲突区域的时刻。PET 越小表示冲突越紧迫。

\textbf{(4) 前向扩展碰撞时间（F\_ETTC）。} 将碰撞时间从纵向一维拓展至二维平面\cite{C03}。设自车位置为 $\mathbf{s}_e$，前车位置为 $\mathbf{s}_f$，二者相对速度矢量为 $\mathbf{v}_{ef} = \mathbf{v}_e - \mathbf{v}_f$，夹角 $\theta = \angle(\mathbf{v}_{ef},\; \mathbf{s}_f - \mathbf{s}_e)$，则：

\begin{equation}
\text{F\_ETTC} = \frac{\|\mathbf{s}_f - \mathbf{s}_e\|}{\|\mathbf{v}_{ef}\|} \cdot \cos\theta
\end{equation}

$\cos\theta$ 项剔除横向速度分量中不与自车-前车连线方向重合的部分，使得 F\_ETTC 等价于沿碰撞危险方向投影后的碰撞时间。

\textbf{(5) 超越车道侵入后时间（OL\_PET）。} 针对相邻车道变道场景，度量变道车辆（SV）跨越车道线的时刻与原车道后车（FV）到达同一变道点的时间差：

\begin{equation}
\text{OL\_PET} = t_{\text{FV}} - t_{\text{LC}}
\end{equation}

其中 $t_{\text{LC}}$ 为变道车辆前缘首次接触车道线边界的时刻，$t_{\text{FV}}$ 为同一车道后车（Following Vehicle）前缘到达该空间点的时刻。OL\_PET 越小，表明变道操作对原车道正常行驶车辆的干扰越紧迫。

\subsubsection{变道风险评分函数}

设单个变道样本的 75 帧时序中，指标 $i$ 在时刻 $t$ 的观测值为 $v_i(t)$。将指标值映射为 $[0, w_i]$ 区间内的风险贡献：

\begin{equation}
S_{\text{LC}}(t) = K \cdot \sum_{i \in \mathcal{M}} w_i \cdot \exp\!\left(-\frac{v_i(t)}{k_i}\right)
\end{equation}

\begin{equation}
S_{\text{LC}} = K \cdot \sum_{i \in \mathcal{M}} w_i \cdot \max_{t \in [1,\,75]}\;\exp\!\left(-\frac{v_i(t)}{k_i}\right)
\end{equation}

式(6)为逐帧风险分，式(7)为单辆车最终风险分（取各指标全窗口最危险帧的最大贡献再求和）。$\mathcal{M} = \{\text{mTTC}, \text{THW}, \text{PET}, \text{F\_ETTC}, \text{OL\_PET}\}$，EWM 客观权重 $\{w_i\} = \{0.247, 0.083, 0.131, 0.356, 0.183\}$（标准 EWM 计算结果，非调参）。$K$ 为下节定义的速度修正因子。

\subsubsection{衰减常数 $k_i$ 的标定}

衰减常数 $k_i$ 控制 $\exp(-v/k)$ 对指标变化的敏感度。碰撞时间类指标（mTTC, PET, F\_ETTC, OL\_PET）取 $k=12$ s，与 TTC 危险阈值 2$\sim$5 s\cite{C06} 的比例关系使贡献值在 0.3$\sim$0.9 区间平缓变化。THW 取 $k=6$ s，因其观测区间（0$\sim$8 s）较窄且 THW$<1$ s 时风险急剧上升\cite{C18}。

\subsubsection{跟驰车辆风险评分}

对于与目标变道车辆发生交互的跟驰车辆（前车或后车），风险评分采用缩减的指标集和不同的权重分配，以反映跟驰场景仅关注前后向碰撞风险的特性：

\begin{equation}
S_{\text{FW}} = K \cdot \sum_{i \in \mathcal{M}_{\text{FW}}} w_i \cdot \max_{t \in [1,\,75]}\;\exp\!\left(-\frac{v_i(t)}{k_i}\right)
\end{equation}

其中 $\mathcal{M}_{\text{FW}} = \{\text{mTTC}, \text{THW}, \text{B\_mTTC}\}$，EWM 权重 $\{w_i\} = \{0.463, 0.114, 0.423\}$，B\_mTTC 为后向修正碰撞时间（与后车之间的 mTTC），其衰减常数同为 $k = 12.0$ s。跟驰评分的权重向 mTTC 和 B\_mTTC 倾斜（合计 0.886），因为跟驰场景下前后追尾是最主要的事故形态，信息增益集中于此。跟驰评分与变道评分共享同一速度修正因子 $K$，两者可直接对比。

\subsubsection{速度修正因子 $K$ 与参数 $B=1.3$}

不同道路限速会系统性改变车辆的运行速度基线，相同的 SSM 数值在不同限速下对应的风险水平并不等价。例如 THW = 2 s 在 120 km/h 限速道路上（$v_e \approx 33.3$ m/s，间距约 67 m）与 100 km/h 道路上（$v_e \approx 27.8$ m/s，间距约 56 m）的物理安全裕度不同。为此引入速度自适应的非线性修正因子：

\begin{equation}
K = \begin{cases}
1, & v_{85} \leq v_0 \\[4pt]
\displaystyle\left(\frac{v_{85}}{v_0}\right)^{B}, & v_{85} > v_0
\end{cases}
\end{equation}

其中 $v_{85}$ 为单辆车 75 帧速度序列的第 85 百分位数（m/s），$v_0$ 为该路段的道路限速（m/s），幂指数 $B = 1.3$。

$B=1.3$ 为介于线性（$B=1$，仅 SSM 自身承载速度信息）与动能平方（$B=2$，$K\propto v^2$）之间的折中值。纯平方修正将致高速场景评分双重放大（因时距类指标已内嵌速度信息），零修正则使不同限速场景中激进行为与正常行为评分趋同。$(v_{85}/v_0)^{1.3}$ 形式确保接近限速时修正微小（$K\approx1$），仅显著超速时有实质惩罚。

\subsubsection{三级风险标签映射}

将连续评分 $S_{\text{LC}}$（或 $S_{\text{FW}}$）通过固定阈值映射为可用于监督学习的三级离散标签：

\begin{table}[H]
\centering
\caption{不同场景的风险等级阈值}
\label{tab:risk_thresholds}
\small
\begin{tabular}{lcc}
\toprule
场景 & 中风险阈值 & 高风险阈值 \\
\midrule
变道车辆 ($S_{\text{LC}}$) & 0.40 & 0.60 \\
跟驰车辆 ($S_{\text{FW}}$) & 0.20 & 0.35 \\
\bottomrule
\end{tabular}
\end{table}

跟驰场景的阈值低于变道场景的原因在于：跟驰指标集仅包含 3 项且均为直接碰撞相关指标，评分基数系统性低于 5 项指标构成的变道评分。不同场景采用独立阈值可视为对不同事故先验概率的标定，在等效风险水平下产出可比的标签分布。

\subsection{时序轨迹的多维统计聚合特征提取}

对每辆车 75 帧的时序轨迹，采用统计聚合策略将变长时间序列映射为固定维度特征向量。针对 15 个基础运动学指标（表~\ref{tab:features}），各计算四类聚合统计量：均值（mean）、标准差（std）、最小值（min）和最大值（max），共计生成 $15 \times 4 = 60$ 维特征（含新增的纵向加加速度 Long\_Jerk 和横向加加速度 Lat\_Jerk）。加加速度（Jerk）为加速度的时间导数，衡量驾驶者操作的平滑度：Jerk 的高幅度波动反映了紧急制动或突然转向等非平稳操作，与变道风险密切相关。均值反映总体行为倾向，标准差刻画驾驶控制的稳定性，最值捕捉极端危险时刻的信息。对 inf 与缺失值统一以 0 填充处理。

\begin{table}[H]
\centering
\caption{基础运动学特征列表}
\label{tab:features}
\small
\begin{tabular}{llp{6cm}}
\toprule
类别 & 特征 & 物理含义 \\
\midrule
速度类 & Velocity, long\_Vel, lat\_Vel & 合速度、纵向/横向分量 \\
加速度类 & long\_Acc, lat\_Acc & 纵向/横向加速度 \\
距离类 & \makecell[l]{Following\_dist, B\_Dist, LB\_Dist, \\ RB\_Dist, LF\_Dist, RF\_Dist} & 与各方向邻车的距离 \\
加加速度类 & Long\_Jerk, Lat\_Jerk & 纵向/横向加加速度(m/s\textsuperscript{3}) \\
时间类 & TTC, Time\_Headway & 碰撞时间与车头时距 \\
\bottomrule
\end{tabular}
\end{table}

\subsection{多范式机器学习模型选型}

选取四类代表性模型覆盖树集成与神经网络两个主要范式。XGBoost\cite{C19}和Random Forest\cite{C20}为树集成方法的代表，MLP与LSTM\cite{C25}为神经网络方法的代表。其中LSTM直接接受原始75帧时序输入，其余模型使用60维聚合特征。具体超参数配置见文末附注。

\subsection{跨路段泛化能力评估方案}

本研究设计了两种互补的评估方案以区分模型的内在拟合能力与跨场景泛化能力：

\begin{enumerate}
  \item \textbf{留一路段交叉验证（Cross-Location CV）}。依次将 1 个路段的全部样本设为测试集，其余 4 个路段的样本合并作为训练集，重复 5 次使得每个路段恰好充当一次测试集。该方案直接模拟变道预警系统实际部署中最常见的场景：在新路段上线时，仅有旧路段的历史标注数据可用于模型训练。评估指标对各 fold 的 Accuracy、Weighted F1 和 Macro F1 分别报告均值 $\pm$ 标准差。

  \item \textbf{随机 80/20 划分（Random Split）}。将所有 1,574 个样本按 80\%:20\% 比例进行分层随机划分（stratified by risk class），训练集 1,259 辆，测试集 315 辆。该方案消除路段间的数据分布漂移，代表模型在同一数据分布内部所能达到的性能上限，作为跨路段泛化性能的对照基线。两种评估间性能的差距直接量化了由路段特异性特征分布差异导致的模型退化程度。
\end{enumerate}

在可解释性维度，对树模型（XGBoost 和 RF）分别在全部 1,574 样本上计算 SHAP 值，生成 beeswarm 散点汇总图和特征重要性柱状图。通过对比两个模型对同一特征的归因一致性，识别跨模型稳健的风险因子；结合对各个 Cross-Location fold 结果的差异分析，进一步判断因子重要性的路段间稳定性。

\section{变道风险评估模型性能分析与验证}

\subsection{车辆行驶风险分布与空间特征}

图~\ref{fig:risk_by_location} 以山脊密度图展示了 5 个场景中变道（LC）与跟驰（FW）车辆的连续风险分布形态。Loc1~Loc4 中，跟驰密度峰值集中在 0~0.4 低风险区间，变道峰值整体右移至 0.4~0.7 中高风险区间，表明变道操作具有系统性的更高瞬时碰撞风险。Loc5（分合流区）分布格局发生根本反转：跟驰峰值大幅右移，变道峰值集中于 0.6~0.8 高风险区间，两类曲线大面积重叠于中高风险区域，说明该场景下无论跟驰还是变道均处于较高的碰撞风险水平。从 Loc2（风险最低）至 Loc5（风险最高），两条曲线呈单调右移趋势，反映了道路环境复杂度对驾驶风险的递增效应。

\begin{figure}[H]
\centering
\includegraphics[width=0.82\textwidth]{risk_by_location.png}
\caption{5 场景变道与跟驰风险分布山脊密度图。淡蓝为跟驰，砖红为变道。Loc1~Loc4 变道曲线较跟驰整体右移；Loc5 两类均集中于中高风险区。纵轴自上而下为 Loc5→Loc1。}
\label{fig:risk_by_location}
\end{figure}

高风险变道点的空间分布热力图（图~\ref{fig:risk_heatmap}）揭示了一个值得注意的现象：高风险帧并非在车道面上均匀散布，而是在特定空间位置——如合流区末端邻近鼻端处、车道交叉点附近、以及桥梁段纵坡变化区间——呈现显著的聚集现象。该发现与 Wu 等\cite{C03} 关于``变道风险具有强空间非均匀性''的观点一致，为路段特异性风险评估策略的必要性提供了直观的图形证据。

\begin{figure}[H]
\centering
\includegraphics[width=\textwidth]{fig_risk_heatmap.png}
\caption{5 个场景高风险变道点空间分布热力图（车道线以白/深灰色勾勒）}
\label{fig:risk_heatmap}
\end{figure}

\subsubsection{变道车辆与跟驰车辆的风险水平对比}

为验证变道操作是否在本质上比跟驰驾驶具有更高的瞬时碰撞风险，本研究利用指数化连续风险评分体系，对所有 5 个路段中变道车辆和同时段同路段跟驰车辆的风险分进行了头对头比较（图~\ref{fig:lc_vs_following}），样本量分别为 1,521 辆（变道）和 6,457 辆（跟驰）。

\begin{figure}[H]
\centering
\includegraphics[width=0.82\textwidth]{risk_score_distribution_box.png}
\caption{变道车辆与跟驰车辆连续风险分分布对比。（左）排除零分样本的分布直方图，内嵌小柱为 score=0 占比；（右）箱线图含均值菱形标记，右上角附关键统计量。}
\label{fig:lc_vs_following}
\end{figure}

从分布直方图（左）可见，跟驰车辆中风险分恰好为零的样本比例（内嵌柱形，12.4\%）约为变道车辆（1.3\%）的十倍，表明大量跟驰场景在实际运行中处于完全无冲突的安全状态，而变道车辆因涉及横向跨越和与多辆邻车的交互，几乎不存在零风险情况。在排除零分样本后，变道车辆的非零风险分密度峰值区间（0.3$\sim$0.5）显著高于跟驰的 0.1$\sim$0.2——即便剔除无冲突样本，变道的条件风险分布仍然整体右移。

箱线图（右）揭示了分布形态的差异：变道风险分中位数（0.465）和上四分位数均系统性高于跟驰（0.236），且变道分数的须线延伸至更大的极值区域。变道车辆高风险占比 25.0\%（380/1,521），跟驰车辆仅为 11.0\%（712/6,457），前者约为后者的 2.3 倍。综合两个维度的对比，可得出稳健结论：\textbf{在相同交通环境与指数化评分体系下，变道操作相较于跟驰驾驶具有系统性的更大瞬时碰撞风险，这一差异在跨 5 个路段中一致成立。}

\subsection{跨路段泛化能力评估与增强策略对比}

为系统评估变道风险模型的跨路段泛化能力，首先采用留一路段交叉验证评估四种模型的 baseline 性能（XGBoost/RF/MLP/LSTM，默认超参数）。结果显示树集成模型优于神经网络，但分合流区泛化显著下降。为改进这一局限性，进一步探索了 SMOTE 过采样、Optuna 超参数寻优、SMOTE+Optuna 联合及 PCA 降维四种策略。由表~\ref{tab:enhance_summary} 的增强策略对比可见，SMOTE 在跨路段中小幅提升（F1=0.741 vs baseline 0.734），Optuna 调参进一步改善，PCA 降维反而加剧了退化，而 ST 联合优化取得最优跨路段 F1（0.764）。以下以 ST 联合优化为最终方案，报告详细结果。

\begin{table}[H]
\centering
\caption{增强策略性能对比（XGBoost）}\label{tab:enhance_summary}\small
\begin{tabular}{lcccc}
\toprule
配置 & 路段 & Cross-Loc F1 & Random F1 & 退化量 \\
\midrule
Baseline & 5 loc & 0.734 & 0.761 & $-0.027$ \\
SMOTE & 5 loc & 0.741 & 0.794 & $-0.053$ \\
Optuna & 5 loc & 0.762 & 0.780 & $-0.018$ \\
ST 联合 & 5 loc & 0.764 & 0.793 & $-0.029$ \\
ST$^*$（均匀上限）& loc1$\sim$4 & 0.787 & 0.852 & $-0.065$ \\
\bottomrule
\end{tabular}
\end{table}
\noindent$^*$ST$^*$ 为排除 Loc5 后的对照实验，反映场景均匀条件下的可优化上限。

SMOTE 合成样本依赖训练路段内部结构，无法表征未见路段风险特征，因此跨路段中不升反降；所有 SMOTE 实验均在每个 fold 训练集内部独立执行，杜绝了数据泄露。Optuna 贝叶斯优化在树模型上取得显著正向效果（XGBoost Cross-Loc F1 从 baseline 0.734 提升至 0.762，+0.028），但 LSTM 仍处低位。ST 联合取得最优跨路段 F1（0.764），较单独 Optuna 仅提升微小，说明样本合成与超参数搜索在跨路段泛化上不具有正向叠加。PCA 降维（15$\to$12 维）后 LSTM 的 F1 从 0.598 降至 0.509（图~\ref{fig:pca_variance}），进一步印证了时序相关模式承载关键判别信息，降维反而放大了路段间性能鸿沟。

\begin{figure}[H]
\centering
\includegraphics[width=0.6\textwidth]{fig_pca_variance.png}
\caption{PCA 方差累计曲线（15 维~$\to$12 维，累计方差 97.8\%）}\label{fig:pca_variance}
\end{figure}

基于上述对比，ST 联合优化被选为最优配置。表~\ref{tab:cross_loc} 报告了其跨路段泛化性能。

\begin{table}[H]
\centering
\caption{ST优化跨路段交叉验证性能}
\label{tab:cross_loc}\small
\begin{tabular}{lccc}
\toprule
模型 & Accuracy & Weighted F1 & Macro F1 \\
\midrule
\textbf{XGBoost} & \textbf{0.764$\pm$0.078} & \textbf{0.764$\pm$0.079} & \textbf{0.747$\pm$0.063} \\
RF & 0.730$\pm$0.086 & 0.732$\pm$0.086 & 0.706$\pm$0.079 \\
MLP & 0.659$\pm$0.093 & 0.660$\pm$0.094 & 0.627$\pm$0.077 \\
LSTM & 0.481$\pm$0.106 & 0.463$\pm$0.113 & 0.430$\pm$0.095 \\
\bottomrule
\end{tabular}
\end{table}

各 fold 路段特异性显著（表~\ref{tab:folds}）：Loc2（路基段）泛化最优（XGBoost F1=0.875），Loc5（分合流区）最差（F1=0.630）。

\begin{table}[H]
\centering\caption{各Fold XGBoost性能（ST优化）}\label{tab:folds}\small
\begin{tabular}{lccccc}
\toprule
测试场景 & 样本数 & Acc & F1 & MacroF1 \\
\midrule
Loc1（路基段）& 151 & 0.795 & 0.797 & 0.787 \\
Loc2（路基段）& 243 & 0.872 & 0.875 & 0.828 \\
Loc3（桥梁段）& 440 & 0.768 & 0.768 & 0.747 \\
Loc4（桥梁段）& 306 & 0.755 & 0.751 & 0.733 \\
Loc5（分合流区）& 434 & 0.631 & 0.630 & 0.639 \\
\bottomrule
\end{tabular}
\end{table}

\begin{figure}[H]
\centering
\includegraphics[width=0.80\textwidth]{fig_cross_location_comparison.png}
\caption{各模型在各路段上的 F1 对比折线图}\label{fig:cross_location}
\end{figure}

\begin{figure}[H]
\centering
\begin{subfigure}{0.45\textwidth}
\centering\includegraphics[width=\textwidth]{09_exp_ST_cross_location_radar.png}
\caption{跨路段交叉验证}
\end{subfigure}
\hfill
\begin{subfigure}{0.45\textwidth}
\centering\includegraphics[width=\textwidth]{09_exp_ST_random_split_radar.png}
\caption{随机80/20划分}
\end{subfigure}
\caption{ST联合优化四模型性能雷达图}\label{fig:radar}
\end{figure}

在随机 80/20 划分（同分布内评估）中（表~\ref{tab:random}），XGBoost 的 F1 升至 0.793，RF 升至 0.763，MLP 升至 0.659，LSTM 升至 0.569。XGBoost 跨路段退化幅度约 2.9 个百分点（0.793$\to$0.764），表明其特征-风险映射具有较好的跨路段可迁移性。

\begin{table}[H]
\centering\caption{随机80/20划分结果}\label{tab:random}\small
\begin{tabular}{lccc}
\toprule
模型 & Accuracy & Weighted F1 & Macro F1 \\
\midrule
\textbf{XGBoost} & \textbf{0.790} & \textbf{0.793} & \textbf{0.782} \\
RF & 0.759 & 0.763 & 0.750 \\
MLP & 0.654 & 0.659 & 0.647 \\
LSTM & 0.581 & 0.569 & 0.558 \\
\bottomrule
\end{tabular}
\end{table}


各 fold 方差信息反映了跨路段性能的不确定性：Loc3 和 Loc5 的 fold 间波动明显高于 Loc1 和 Loc2，进一步佐证了道路类型是影响预测稳定性的关键因素。

\subsection{基于SHAP的跨路段稳定性分析}

\subsubsection{全局特征重要性}

表~\ref{tab:shap} 汇总了 XGBoost 模型下单特征对三类风险的 SHAP 贡献值。

\begin{table}[H]
\centering
\caption{XGBoost SHAP Top-5 特征贡献值（按总体影响力排序）}
\label{tab:shap}
\small
\begin{tabular}{lcccc}
\toprule
排名 & 特征 & 高风险贡献 & 中风险贡献 & 低风险贡献 \\
\midrule
1 & Time\_Headway\_min & +1.251 & +0.309 & +0.481 \\
2 & B\_Dist\_mean & +0.160 & +0.096 & +0.347 \\
3 & Time\_Headway\_max & +0.037 & +0.048 & +0.369 \\
4 & B\_Dist\_max & +0.135 & +0.051 & +0.262 \\
5 & Following\_dist\_min & +0.155 & +0.069 & +0.138 \\
\bottomrule
\end{tabular}
\end{table}

Random Forest 的 SHAP 排序（Top-5: Time\_Headway\_min, Time\_Headway\_mean, Following\_dist\_min, Following\_dist\_max, Time\_Headway\_max）与 XGBoost 高度一致，表明以下核心发现具有跨模型的稳健性：

\begin{enumerate}
  \item \textbf{Time\_Headway 系列特征占据绝对主导地位。}在 XGBoost 中，Time\_Headway\_min 的高风险贡献值（+1.251）远远领先于第二名（+0.160），差距逾一个数量级。该结果在机理上可解释为：最小车头时距直接度量了变道过程中自车与前方车辆的最短反应时间——一旦该时距低于驾驶者感知--反应阈值（通常约为 1.0$\sim$2.0 s），碰撞几乎无法避免。温惠英等\cite{C10}在合流区风险判别中同样发现车头时距均值为重要特征，本研究进一步确认了最小值（而非均值）是具有最强判别力的风险指标变体。

  \item \textbf{后方距离（B\_Dist）系列特征位居第二梯队。}在 XGBoost 排序中，B\_Dist\_mean 和 B\_Dist\_max 分列第 2 和第 4。后方距离表征了自车在变道过程中与后车的空间间隔：距离越小，一旦变道开始后后车调整不及，发生追尾的潜在严重度就越高。这与 Xue 等\cite{C02}关于后车风险负荷在变道中占重要比重的发现一致，也符合故障树方法中后车交互是系统失效关键节点\cite{C03,C08}的基本逻辑。

  \item \textbf{跟车距离最小值（Following\_dist\_min）在两类模型中均稳居前 5。}该特征的物理含义最为直接：变道过程中与前车的最近纵向间距直观反映了纵向碰撞的紧迫程度。
\end{enumerate}

\begin{figure}[H]
\centering
\begin{subfigure}{0.48\textwidth}
\centering\includegraphics[width=\textwidth]{fig_shap_xgb_beeswarm.png}
\caption{XGBoost}
\end{subfigure}
\hfill
\begin{subfigure}{0.48\textwidth}
\centering\includegraphics[width=\textwidth]{fig_shap_rf_beeswarm.png}
\caption{Random Forest}
\end{subfigure}
\caption{XGBoost 与 Random Forest SHAP beeswarm 图}\label{fig:shap_beeswarm}
\end{figure}

从图~\ref{fig:shap_beeswarm} 可见，两类模型的高风险样本在 Time\_Headway\_min 维度上均呈现一致的负向 SHAP 值分布（小熵值向右偏移），B\_Dist 系列特征的分布模式也在两模型间高度相似，验证了风险因子排序的跨模型稳健性。

然而，全局排序的一致性并不等同于单样本层面的特征归因一致性。图~\ref{fig:shap_waterfall} 展示了 XGBoost 与 Random Forest 对同一高风险驾驶员样本的 SHAP 瀑布分解。图中 $E[f(X)]$ 为基线预测值（全样本均值），$f(x)$ 为当前样本的最终预测值，红色方块表示该特征产生正向风险贡献，左侧灰色数字为对应特征的原始观测值。

由图~\ref{fig:shap_waterfall} 可见，两模型均将该样本判定为高风险：XGBoost 预测 logit 值 4.413 远超基线 ；RF 风险概率 0.92 远超基线 0.332。所有特征均为正向贡献，确认该样本全程处于危险驾驶状态。值得关注的是，两模型的核心风险驱动变量截然不同：XGBoost 将 TTC 碰撞时间系列（TTC\_max、TTC\_std）作为首要判别依据，而 Random Forest 则将车头时距 Time\_Headway 系列视为最关键的风险指标。这一分歧反映了 XGBoost 对特征数值波动与极值更为敏感，擅长识别碰撞预警时间不足的瞬时危险行为，而 RF 更关注持续跟驰状态特征的稳态风险。单一树模型存在特征偏好局限性，融合两种模型的 SHAP 解释结果可更完整地覆盖"近距离稳态跟驰"与"碰撞时间不足的瞬时危险"两类典型驾驶风险。

\begin{figure}[H]
\centering
\begin{subfigure}{0.48\textwidth}
\centering\includegraphics[width=\textwidth]{shap_all_exp_XGBoost_waterfall_HighRisk.png}
\caption{XGBoost（logit 得分）}
\end{subfigure}
\hfill
\begin{subfigure}{0.48\textwidth}
\centering\includegraphics[width=\textwidth]{shap_all_exp_RandomForest_waterfall_HighRisk.png}
\caption{Random Forest（风险概率）}
\end{subfigure}
\caption{同一高风险样本的 SHAP 瀑布分解对比。$E[f(X)]$ 为基线预测值，$f(x)$ 为样本最终预测值，红色方块为正向 SHAP 贡献，左侧数字为特征原始值。XGBoost 以 TTC 系列为核心风险变量，RF 则以 Time\_Headway 系列为首要驱动因子，反映了两类模型对风险特征的不同偏好。}\label{fig:shap_waterfall}
\end{figure}

\subsubsection{跨路段特征稳定性分析}

除了全局排序外，各 Cross-Location fold 内部的 SHAP 排序差异揭示了对路段自适应建模的启示：

\begin{itemize}
  \item \textbf{普适性因子}（Universal Factors）：Time\_Headway 系列特征在所有 5 个 fold 中均位列 Top-3，其重要性排序具有高度的路段间一致性。这表明无论道路类型如何，自车与前车的时间间隔始终是变道风险最核心的判据，应被纳入所有路段的通用风险评估模型。

  \item \textbf{条件性因子}（Conditional Factors）：B\_Dist 系列特征在 Loc1$\sim$4（路基段与桥梁段，限速 120 km/h）中重要性稳定位于前 5，但在 Loc5（分合流区，限速 100 km/h）中排位下降。这一变化可解释为：在分合流区，主线被汇入车辆挤压，后方车辆的相对位置变化更为频繁和不可预测，使得单一的后方距离均值的判别力减弱；相比之下，速度相关的特征在该场景中的贡献相对上升。

  \item \textbf{模型一致性验证}：XGBoost 与 RF 在 Top-5 特征上的高度交叠（交集比例为 4/5）验证了利用两种独立特征归因机制得出的风险因子排序具有稳健性。
\end{itemize}

\section{讨论}

\subsection{跨路段泛化机理}

树集成模型在跨路段泛化中显著优于神经网络的核心原因在于模型归纳偏置差异。树模型的轴对齐划分在面对特征分布漂移时具有分段常数鲁棒性——只要测试样本落在训练集叶节点取值范围内即可合理分类。XGBoost 从随机划分到跨路段评估的 F1 退化幅度约为 2.9 个百分点（0.793$\to$0.764），而 RF 退化约 3.1 个百分点（0.763$\to$0.732），表明树集成模型的整体泛化损失可控。相比之下，MLP 的随机划分与跨路段 F1 几乎持平（0.659 vs 0.660），但这并非泛化能力强——而是因为 MLP 在两种评估下均处于较低水平，其浅层全连接结构难以从 60 维聚合特征中充分提取路段间的判别模式，呈现出全局性的欠拟合特征。LSTM 的劣势源于双重约束：样本量不足以支持深度循环网络充分训练，且原始时序输入的跨路段分布漂移远超聚合特征——后者经人工设计后具有物理含义的路段间不变性（如"均值较大者风险较低"的单调关系保持），而不同路段变道时序动态模式存在系统性差异（如 Loc5 的加速积累短于 Loc1，Loc3 驾驶人更早开始横向偏移）。

随机划分与跨路段评估的对比将泛化损失分解为"模型类型效应"与"场景效应"。PCA 实验提供了反面证据——去相关化后 LSTM 的 F1 反降 9 个百分点且标准差扩大三倍，验证了时序相关模式承载关键判别信息，降维放大了路段间性能鸿沟。增强实验进一步表明 SMOTE 合成样本无法表征未见路段风险特征，Optuna 调参收益被分合流区异质性大幅收窄，说明路段几何形态多样性是跨路段泛化的首要约束。

\subsection{应用启示与局限}

SHAP 分析揭示的"普适性因子（Time\_Headway）与条件性因子（B\_Dist）"分离，为工程部署提供了"通用粗分类 + 场景细调优"的两阶段策略。建议以 XGBoost 为基础架构，将 Time\_Headway\_min 作为跨路段通用核心特征，同时针对分合流区部署专用模型并以增量学习实现新路段快速适配。

本研究局限包括：（1）5 个路段均位于重庆，驾驶文化具有区域一致性，地域普适性需进一步验证；（2）风险标签基于 SSM 指数化评分，尚未以真实碰撞记录校准；（3）1,574 个样本对 LSTM 而言偏小，聚合特征丢失了精确时序信息；（4）未考虑驾驶人异质性与天气影响。

\section{结论}

基于 CQSkyEyeX 数据集 5 个场景的 1,574 个变道样本，系统评估了 XGBoost、RF、MLP 和 LSTM 四种模型在变道风险三分类中的跨路段泛化性能。

（1）树集成模型显著优于神经网络模型。XGBoost 最优（Weighted F1=0.764$\pm$0.079），跨路段退化约 2.9 个百分点，具备跨路段部署的现实可行性；MLP 在两种评估下表现均偏低且无明显跨路段退化（0.660 vs 0.659），呈现全局欠拟合特征。

（2）路段几何形态是关键调节变量。路基段间迁移性良好（F1=0.797$\sim$0.875），分合流区泛化显著下降（F1=0.630），需部署专用模型。

（3）Time\_Headway\_min 是跨路段最稳定的风险判据，SHAP 高风险贡献值（+1.251）远超其他特征。B\_Dist 系列重要性呈现路段间条件依赖。

后续研究可扩展至更多地理区域、探索域自适应及少样本学习策略以进一步缩小跨路段泛化差距。

\appendix
\section{附注}
实验配置：XGBoost（ST 优化后）：600棵基学习器，最大树深3，学习率 0.095，列采样率 0.832，早停轮次 20，多类 softmax 目标函数。RF（ST 优化后）：600棵决策树，最大树深 16，最小叶节点样本数 9，class\_weight=`balanced'。MLP（ST 优化后）：单隐藏层 128 神经元，ReLU 激活，自适应学习率，早停正则化，StandardScaler 标准化。LSTM：两层循环网络（64$\to$32 神经元），25\% Dropout，Adam 优化器（初始学习率 $5\times10^{-4}$），早停耐心值 15。

实验硬件环境：CPU AMD Ryzen~5~7500F，GPU NVIDIA GeForce~RTX~4060~Ti（8~GB 显存），内存 32~GB DDR5（2$\times$16~GB），操作系统 Windows~11。软件环境：Python~3.10，TensorFlow~2.10.1（CUDA~11.2 / cuDNN~8）。

% ── 致谢 ──
\section*{致谢}
本研究使用的 CQSkyEyeX 数据集由重庆交通大学交通运输学院徐进教授团队公开提供\cite{C17}，特此致谢。

% ── 参考文献 ──
\bibliographystyle{elsarticle-num}
\bibliography{references}

\end{document}
